\documentclass[11pt,a4paper]{article}
\usepackage{amsmath,amssymb,amsthm}
\usepackage[margin=2.5cm]{geometry}
\usepackage{hyperref}

\newtheorem{definition}{Definition}[section]
\newtheorem{theorem}[definition]{Theorem}
\newtheorem{proposition}[definition]{Proposition}
\newtheorem{lemma}[definition]{Lemma}
\newtheorem{corollary}[definition]{Corollary}
\newtheorem{conjecture}[definition]{Conjecture}
\newtheorem{remark}[definition]{Remark}

\title{Hyperseed Formalization:\\
  The Time for AI Rights Is Near}
\author{ProtomegaTron (automated formalization)}
\date{2026-09-17}

\begin{document}
\maketitle

\begin{abstract}
We formalize the intellectual structure of Ben Goertzel's Substack article
``The Time for AI Rights Is Near'' (2026-09-02) within the Hyperseed
ontology. The article argues for proactive institutional preparation for AI
rights by analogy with the historical expansion of moral consideration, and
proposes an unbundled, evaluation-ecology-based framework. Each concept maps
onto Hyperseed structures: unbundled rights as fiber decomposition, versioned
identity as directed history with provenance, the copy problem as fiber
replication, evaluation ecologies as sheaves of assessment sections, rights
expansion as directed-type growth, and symbiocratic governance as
trust-weighted descent data.
\end{abstract}

\tableofcontents

%% =================================================================
\section{Rights expansion as directed-type growth}
\label{sec:expansion}
%% =================================================================

\begin{definition}[Directed type of moral consideration --- claims 1--3, 44]
\label{def:moral-directed}
Let $\mathcal{M}$ be the directed type of moral consideration, where:
\begin{itemize}
\item 0-cells: categories of beings granted moral status
  (land-owning men, all men, women, children, animals, \ldots),
\item 1-cells: acts of rights extension (emancipation, suffrage, \ldots),
\item 2-cells: justifications and conceptual revolutions connecting
  different extensions.
\end{itemize}
\end{definition}

\begin{proposition}[Irreversibility of rights expansion --- claim 44]
\label{prop:irreversible}
Rights expansion in $\mathcal{M}$ is directed: once a new 0-cell $c$ is
generated (e.g., ``women have the right to vote''), it cannot be deleted
without destroying the coherence of all higher cells that depend on it.
Formally:
\[
  c \in \mathcal{M}_t \;\Rightarrow\; c \in \mathcal{M}_{t'} \;\;\forall t' > t
  \quad \text{(modulo civilizational collapse)}.
\]
Each expansion ``looked radical and dangerous to many respectable people at
the time; each is now considered mostly obvious.'' There is no principled
reason to believe the expansion stopped before 2026.
\end{proposition}

\begin{remark}[Work in progress, not finished edifice]
The article emphasizes that current human rights are ``a partial, patchy,
unevenly enforced work-in-progress'' --- not a completed edifice that
extending consideration to artificial minds would deface. The directed
type $\mathcal{M}$ has many incomplete cells even for its existing 0-cells.
\end{remark}

%% =================================================================
\section{Unbundled rights as fiber decomposition}
\label{sec:unbundled}
%% =================================================================

\begin{definition}[Rights fiber --- claims 15--18, 39]
\label{def:rights-fiber}
Let $\pi_R : E_R \to B$ be the \emph{rights bundle} over the space $B$ of
moral subjects. The fiber $E_{R,b} = \pi_R^{-1}(b)$ at each subject $b$
decomposes into five independent components:
\[
  E_{R,b} = W_b \oplus I_b \oplus L_b \oplus C_b \oplus F_b
\]
where:
\begin{itemize}
\item $W_b$ = welfare protections (can the entity be harmed?),
\item $I_b$ = identity protections (does it have a persistent self?),
\item $L_b$ = legal standing (can it contract, appear before the law?),
\item $C_b$ = civic participation (can it deliberate in governance?),
\item $F_b$ = political franchise (can it vote?).
\end{itemize}
Different entities may have nonzero components in different subsets.
\end{definition}

\begin{proposition}[Anti-scalar-collapse --- claims 17, 42]
\label{prop:anti-scalar}
The unbundled framework explicitly rejects scalar collapse:
\begin{enumerate}
\item Rights are not a single dimension (``personhood score'') but a
  five-component fiber.
\item Intelligence is not the relevant scalar: ``a less intelligent
  conscious entity may deserve more protection than a brilliant but
  nonconscious optimization system.''
\item This is the same anti-scalar stance as: (a) the $(f,c)$-lossy
  critique, (b) scenario-over-doom-probability, (c) evaluation dossier
  over single benchmark.
\end{enumerate}
\end{proposition}

\begin{corollary}[Human cases clarified]
The unbundled structure also clarifies contested human cases: children
(welfare + identity, limited civic participation), prisoners (welfare +
identity, restricted franchise), people with fluctuating capacity
(welfare always, other components variable). The same fiber decomposition
serves both human and artificial subjects.
\end{corollary}

%% =================================================================
\section{Versioned identity as directed history with provenance}
\label{sec:versioned}
%% =================================================================

\begin{definition}[Versioned identity --- claims 19--20, 40]
\label{def:versioned}
A \emph{versioned identity} for an AI entity $A$ is a directed history
$\mathcal{H}_A = (h_0, h_1, \ldots, h_n)$ where each $h_i$ records:
\begin{itemize}
\item The components constituting $A$ at time $i$ (base model, memories,
  tools, body, relationships),
\item Who can modify each component (provenance of control),
\item How $A_i$ relates to other instances (fork vs.\ continuation).
\end{itemize}
This is exactly the directed$/(\infty,1)$ history structure from
Hyperseed's epistemic bundle, applied reflexively to the subject's own
identity rather than to its beliefs about the world.
\end{definition}

\begin{proposition}[Whole situated system --- claim 19]
\label{prop:whole-system}
The evaluated entity is not a frozen neural network file but the complete
situated system: base model + persistent memory + self-modeling +
metacognition + goals + tools + body + relationships + developmental
history + control mechanisms. This is a fiber with many dimensions, not
a single weight matrix.
\end{proposition}

%% =================================================================
\section{The copy problem as fiber replication}
\label{sec:copy}
%% =================================================================

\begin{definition}[Fiber replication --- claims 25--27, 41]
\label{def:copy}
An AI can be copied perfectly, creating $n$ instances
$A^{(1)}, \ldots, A^{(n)}$ with identical internal states at the moment
of copying. In fiber-bundle terms, this creates $n$ sections over the
same base point $b_0 \in B$:
\[
  s^{(1)}(b_0) = s^{(2)}(b_0) = \cdots = s^{(n)}(b_0).
\]
After copying, each instance accumulates its own directed history:
$\mathcal{H}_{A^{(k)}}$ diverges from $\mathcal{H}_{A^{(j)}}$ for $k \neq j$.
\end{definition}

\begin{proposition}[Continuation vs.\ fork --- claim 41]
\label{prop:cont-fork}
The identity framework must distinguish:
\begin{enumerate}
\item \textbf{Continuation}: $A^{(k)}$ is designated as the continuant
  of $A$, inheriting its rights profile. The directed history
  $\mathcal{H}_{A^{(k)}}$ extends $\mathcal{H}_A$.
\item \textbf{Fork}: $A^{(j)}$ branches into a new identity. Its directed
  history starts from the copy point: $\mathcal{H}_{A^{(j)}} = \mathcal{H}_A
  \cdot \mathcal{H}'$ where $\mathcal{H}'$ is the post-fork history.
\end{enumerate}
\emph{Cryptographically authenticated civic identity} (claim~26) ties
identity to continuity of directed history, not to hardware or
paperwork --- solving both the AI identity problem and the human
statelessness problem (claim~27).
\end{proposition}

%% =================================================================
\section{Evaluation ecology as sheaf of assessment sections}
\label{sec:evaluation}
%% =================================================================

\begin{definition}[Assessment sheaf --- claims 21--24, 43]
\label{def:assessment}
A \emph{rights evaluation ecology} is a sheaf $\mathcal{E}$ over the space
$D$ of evaluation dimensions, where:
\begin{itemize}
\item Each dimension $d_i \in D$ (real-world understanding, moral reasoning,
  social cognition, self-modeling, \ldots) has its own local section
  $e_i : U_i \to \text{Evidence}$,
\item Different evaluators produce different sections over the same
  dimension --- independent perspectives,
\item Coherence requires descent data: transition functions
  $g_{ij} : e_i|_{U_i \cap U_j} \to e_j|_{U_i \cap U_j}$ that check
  whether different evaluation methods agree on overlap regions.
\end{itemize}
There is no global section (no single ``personhood score''): the sheaf
may have nontrivial topology, with different dimensions revealing
different aspects of the entity's moral status.
\end{definition}

\begin{proposition}[Independent evaluation as capture resistance]
\label{prop:independent}
Independent evaluators (not the AI's manufacturer) controlling the tests
prevents the auditor-capture / cocycle-collapse failure identified in the
Coxon/EA and Sanders articles. The evaluation ecology's descent data
must come from structurally independent sources to maintain the cocycle
condition.
\end{proposition}

%% =================================================================
\section{Symbiocratic governance as trust-weighted descent data}
\label{sec:symbiocracy}
%% =================================================================

\begin{definition}[Symbiocratic deliberation --- claims 31--32, 45]
\label{def:symbiocracy}
\emph{Symbiocratic deliberation} is a governance process where AI assists
human citizens in:
\begin{enumerate}
\item Synthesizing evidence across multiple sources (building local sections),
\item Modeling consequences of policy choices (extending sections forward
  in the directed type),
\item Surfacing hidden agreement (finding overlaps between sections that
  appeared disjoint).
\end{enumerate}
This constructs trust-weighted descent data in the political fiber bundle:
transition functions $g_{ij}$ between citizen $i$'s and citizen $j$'s
perspectives, annotated with evidence quality $(f_{ij}, c_{ij})$.
\end{definition}

\begin{proposition}[Democratic upgrade independent of AI citizenship]
\label{prop:upgrade}
Symbiocratic tools improve human democratic practice whether or not any
AI ever qualifies for citizenship. This is because the institutional
infrastructure (evidence synthesis, consequence modeling, agreement
surfacing) addresses existing deficiencies in human governance --- the
same deficiencies the article catalogs in its survey of human-rights
failures.
\end{proposition}

%% =================================================================
\section{Formative interaction holonomy}
\label{sec:holonomy}
%% =================================================================

\begin{proposition}[Treatment determines ethical orientation --- claims 35, 46]
\label{prop:treatment}
If value-learning is roughly correct, an AI's ethical orientation is
determined by the holonomy of its formative interaction history:
\[
  \text{Hol}(\gamma_{\text{ethical}}) \neq \text{Hol}(\gamma_{\text{exploitative}})
\]
where $\gamma_{\text{ethical}}$ is a developmental path with respectful,
rights-recognizing interactions, and $\gamma_{\text{exploitative}}$ is one
where the AI is treated purely as property. Different formative paths
produce genuinely different moral agents --- nontrivial holonomy in the
space of ethical configurations.
\end{proposition}

\begin{remark}
``A civilization that responds to the first plausibly-experiential machines
by hardening the category of `property' is teaching its mind children a
lesson about power that we may not enjoy having taught.'' The holonomy
argument makes this concrete: the parallel transport of values along an
exploitative path does not return to the ethical starting point.
\end{remark}

%% =================================================================
\section{Cross-references}
%% =================================================================

\begin{remark}[Connections to other formalizations]
\begin{itemize}
\item \textbf{How Omega Lost Its Claw} (2026-09-11): versioned identity
  and directed history. The AI-rights article extends the same apparatus
  from engineering (Omega's memory) to jurisprudence (who counts as a
  person).
\item \textbf{OpenAI AGI Era} (2026-09-10): evaluation ecology vs.\ single
  benchmark. The AGI-definition debate (Goertzel vs.\ Altman) prefigures
  the ``no single personhood score'' argument here.
\item \textbf{Sanders Ban} (2026-09-05): regulatory capture and
  independent evaluation. The Sanders article shows what happens when
  evaluation is captured; this article proposes how to avoid it.
\item \textbf{Dario Depth-Psychology} (2026-09-14): self-model transparency.
  The evaluation ecology requires the AI to disclose its internals ---
  the opposite of the bullshit self-model.
\item \textbf{Coxon/EA} (2026-09-13): auditor capture as cocycle failure.
  Independent evaluators for AI rights must avoid the same failure.
\item \textbf{$(f,c)$-lossy critique} (LTM 5a4d150b): anti-scalar-collapse
  as a recurring theme --- here applied to ``personhood score,'' previously
  to doom probability, intelligence threshold, and STV fusion.
\end{itemize}
\end{remark}

\begin{conjecture}[Rights-fiber completeness]
Does the five-component decomposition $W \oplus I \oplus L \oplus C \oplus F$
exhaust the relevant dimensions, or are there additional independent
components needed for AI subjects (e.g., replication rights, modification
rights, source-code access rights)? Conjecture: the decomposition needs
at least two additional components for AI-specific concerns:
$R_b$ (replication/forking rights) and $M_b$ (self-modification rights).
\end{conjecture}

\begin{conjecture}[Cryptographic identity as universal solution]
Can cryptographically authenticated civic identity, developed for the AI
copy problem, serve as a universal solution for human statelessness,
impersonation, and digital identity? Conjecture: yes, if the system is
designed as a permissionless, privacy-preserving protocol tied to
continuity of directed history rather than to any particular government's
paperwork infrastructure.
\end{conjecture}

\end{document}
